\section{Hidden Markov models from zero}\label{sec:hmm}

\subsection{A chain of phases}\label{sec:chain}

Number the fixes $k=1,2,\dots,K$ and call $Q_k$ the phase of fix $k$. Suppose for a moment that we
could see it. The simplest model of how phases follow each other is a \emph{Markov chain}: the
phase of the next fix depends only on the phase of the current one,
\begin{equation}\label{eq:markov}
  P(Q_{k+1}=j \mid Q_k=i,\ Q_{k-1},\dots,Q_1) = a_{ij}.
\end{equation}
The numbers $a_{ij}$ form the \emph{transition matrix} $A$. Row $i$ lists where phase $i$ goes
next, so each row sums to 1. The first phase is drawn from the \emph{initial distribution}
$\nu_j=P(Q_1=j)$. Vilpellet's chain, one step per second (Figure~\ref{fig:chain}), is
\begin{equation}\label{eq:A}
  A=\begin{pmatrix} 0.9896 & 0.0035 & 0.0069\\ 0.0067 & 0.9900 & 0.0033\\ 0.0010 & 0.0071 & 0.9919 \end{pmatrix}
  \qquad\text{(rows and columns: \cl, \se, \tr).}
\end{equation}

The diagonal says how sticky each phase is: a climb goes on to the next fix with probability
0.9896. A run of phase $i$ therefore lasts $n$ fixes with probability $a_{ii}^{\,n-1}(1-a_{ii})$, a
geometric law with mean $1/(1-a_{ii})$: 96~s for a climb, 100~s for a search, 123~s for a
transition. The chain has no memory of how long it has been in a phase: a climb that has lasted
ten minutes ends at the next fix with the same probability as one that has just begun.

\begin{figure}
\centering
\begin{tikzpicture}[x=1cm,y=1cm, every loop/.style={min distance=11mm, looseness=6}]
  \node[state=climb] (C) at (0,0) {\cl\\[-1pt]{\footnotesize 96 s}};
  \node[state=search] (S) at (4.2,2.4) {\se\\[-1pt]{\footnotesize 100 s}};
  \node[state=transition] (T) at (8.4,0) {\tr\\[-1pt]{\footnotesize 123 s}};
  \path[->, ink2, line width=0.7pt, font=\footnotesize]
    (C) edge[loop left] node[left] {0.9896} (C)
    (S) edge[loop above] node[above] {0.9900} (S)
    (T) edge[loop right] node[right] {0.9919} (T)
    (C) edge[bend left=12] node[above left=-2pt] {0.0035} (S)
    (S) edge[bend left=12] node[below right=-2pt] {0.0067} (C)
    (C) edge[bend left=10] node[above] {0.0069} (T)
    (T) edge[bend left=10] node[below] {0.0010} (C)
    (S) edge[bend left=12] node[above right=-2pt] {0.0033} (T)
    (T) edge[bend left=12] node[below left=-2pt] {0.0071} (S);
\end{tikzpicture}
\caption{Vilpellet's chain of phases. Arrows carry the probability of the next fix's phase given the
current one; under each name, the mean duration $1/(1-a_{ii})$ seconds that the chain alone implies.}
\label{fig:chain}
\end{figure}

\begin{inwords}
$A$ is a table of habits: what usually follows what, one second at a time. It knows nothing about
the track.
\end{inwords}

\subsection{Hidden phases, visible symbols}\label{sec:hidden}

We never see $Q_k$. At each fix we see something else, $X_k$: in Vilpellet's method one of eight
symbols computed from the track (Section~\ref{sec:method}). Each phase produces the symbols with
its own probabilities, the \emph{emission probabilities}
\begin{equation}\label{eq:emission}
  b_j(x) = P(X_k=x \mid Q_k=j).
\end{equation}
The model assumes that once the phase of fix $k$ is known, its symbol depends on nothing else:
neither on the other phases nor on the other symbols. The probability of a phase sequence
$q_1,\dots,q_K$ together with the symbols $x_1,\dots,x_K$ is then a product of one factor per
fix (Figure~\ref{fig:graphical}):
\begin{equation}\label{eq:joint}
  P(q_{1:K},x_{1:K}) = \nu_{q_1}\,b_{q_1}(x_1)\prod_{k=2}^{K} a_{q_{k-1}q_k}\,b_{q_k}(x_k).
\end{equation}
This is a \emph{hidden Markov model} (HMM)~\cite{rabiner1989,zucchini2016}: a Markov chain that we
cannot see, observed through noisy symbols. Movement ecology labels the behaviour of animals from
GPS tracks with the same model, taking as observations the step length and the turning angle of
each fix, drawn from a law per state such as a gamma and a von
Mises~\cite{langrock2012,michelot2016}. Vilpellet's observations are of another kind: yes/no
summaries of 30 fixes of track, which Section~\ref{sec:assumptions} returns to.

\begin{figure}
\centering
\begin{tikzpicture}[x=1cm,y=1cm,
    hid/.style={circle, draw=ink, line width=0.8pt, minimum size=9mm, font=\small},
    obs/.style={rectangle, draw=ink, fill=light, line width=0.8pt, minimum size=8mm, font=\small}]
  \foreach \k/\x in {1/0, 2/2.2, 3/4.4, 4/6.6} {
    \node[hid] (q\k) at (\x,1.6) {$Q_{\k}$};
    \node[obs] (x\k) at (\x,0) {$X_{\k}$};
    \draw[->, ink2, line width=0.8pt] (q\k) -- node[right, lbl] {$b$} (x\k);
  }
  \node[hid] (qK) at (9.4,1.6) {$Q_K$};
  \node[obs] (xK) at (9.4,0) {$X_K$};
  \draw[->, ink2, line width=0.8pt] (qK) -- node[right, lbl] {$b$} (xK);
  \foreach \a/\b in {q1/q2, q2/q3, q3/q4} \draw[->, ink2, line width=0.8pt] (\a) -- node[above, lbl] {$A$} (\b);
  \draw[->, ink2, line width=0.8pt, dashed] (q4) -- (qK);
  \node[note, anchor=west] at (10.2,1.6) {phases: hidden};
  \node[note, anchor=west] at (10.2,0) {symbols: seen};
\end{tikzpicture}
\caption{A hidden Markov model. Each phase depends only on the previous one, through $A$; each
symbol depends only on the phase of its own fix, through $b$.}
\label{fig:graphical}
\end{figure}

\subsection{A small example}\label{sec:toy}

Take two phases, climb $C$ and glide $G$, and one yes/no symbol: does the pilot rise (1) or sink
(0) at this step? Let both phases start with probability $0.5$, stay with probability $0.9$ and
switch with $0.1$, and let $b_C(1)=0.9$ and $b_G(1)=0.2$. The pilot shows
\[
  x = 1,\ 1,\ 0,\ 1,\ 1,\ 0,\ 0,\ 0 .
\]
Reading each symbol alone, the likelier phase is $C$ after a 1 and $G$ after a 0, which gives
$C\,C\,G\,C\,C\,G\,G\,G$: the single 0 at step 3 makes the pilot glide for one step. With the
chain, compare two sequences that differ only at step 3. By~\eqref{eq:joint} all their other
factors are equal, and
\[
  \frac{P(\dots C\,C\,C\dots)}{P(\dots C\,G\,C\dots)}
  = \frac{a_{CC}\,b_C(0)\,a_{CC}}{a_{CG}\,b_G(0)\,a_{GC}}
  = \frac{0.9\cdot0.1\cdot0.9}{0.1\cdot0.8\cdot0.1}\approx10 .
\]
Staying in the climb costs one unlikely symbol, $0.1$ instead of $0.8$; leaving and coming back
costs two unlikely changes, $0.1\times0.1$. The isolated sink is read as noise. The zeros from
step 6 on persist, and the best sequence switches to $G$ at step 6
(Figure~\ref{fig:trellis}). This is the job of the HMM in Vilpellet's method: a phase changes only
when the evidence persists.

\begin{figure}
\centering
\begin{tikzpicture}[x=1.65cm,y=1cm,
    cell/.style={rectangle, draw=muted, rounded corners=1pt, minimum width=1.0cm, minimum height=0.6cm, font=\footnotesize, fill=white}]
  \foreach \k/\ob/\vc/\vg in {1/1/.818/.182, 2/1/.953/.047, 3/0/.529/.471, 4/1/.835/.165,
                              5/1/.958/.042, 6/0/.529/.471, 7/0/.123/.877, 8/0/.017/.983} {
    \node[lbl] at (\k,2.05) {$k=\k$};
    \node[font=\small\bfseries] at (\k,1.6) {$x=\ob$};
    \node[cell] (c\k) at (\k,0.8) {\vc};
    \node[cell] (g\k) at (\k,-0.4) {\vg};
  }
  \node[lbl, anchor=east] at (0.45,0.8) {$\delta_k(C)$};
  \node[lbl, anchor=east] at (0.45,-0.4) {$\delta_k(G)$};
  \foreach \k/\pc/\pg in {2/0/1, 3/0/0, 4/0/1, 5/0/1, 6/0/0, 7/0/1, 8/0/1} {
    \pgfmathtruncatemacro{\j}{\k-1}
    \ifnum\pc=0 \draw[->, muted] (c\j.east) -- (c\k.west); \else \draw[->, muted] (g\j.east) -- (c\k.west); \fi
    \ifnum\pg=0 \draw[->, muted] (c\j.east) -- (g\k.west); \else \draw[->, muted] (g\j.east) -- (g\k.west); \fi
  }
  \foreach \k in {1,...,5} \node[cell, draw=climb, line width=1.2pt, fill=climb!12] at (c\k) {};
  \foreach \k/\v in {1/.818, 2/.953, 3/.529, 4/.835, 5/.958} \node[font=\footnotesize] at (c\k) {\v};
  \foreach \k in {6,7,8} \node[cell, draw=transition, line width=1.2pt, fill=transition!12] at (g\k) {};
  \foreach \k/\v in {6/.471, 7/.877, 8/.983} \node[font=\footnotesize] at (g\k) {\v};
  \foreach \a/\b in {c1/c2, c2/c3, c3/c4, c4/c5} \draw[->, climb, line width=1.4pt] (\a.east) -- (\b.west);
  \draw[->, ink, line width=1.4pt] (c5.east) -- (g6.west);
  \foreach \a/\b in {g6/g7, g7/g8} \draw[->, transition, line width=1.4pt] (\a.east) -- (\b.west);
\end{tikzpicture}
\caption{Viterbi on the small example. Each cell holds $\delta_k(j)$, the probability of the best
sequence ending in phase $j$ at step $k$, divided by the sum of the column; the grey arrow into each
cell comes from its best predecessor. Steps 3 and 6 have the same scores, yet the path stays in $C$
at step 3 and leaves it at step 6: the decision depends on the symbols that follow, which is why
Viterbi reads the whole sequence before deciding.}
\label{fig:trellis}
\end{figure}

\subsection{Decoding and fitting}\label{sec:three}

\paragraph{The best phase sequence (Viterbi).} Write $\delta_k(j)$ for the highest
probability~\eqref{eq:joint} of any phase sequence of the first $k$ fixes that ends in phase $j$.
The best sequence ending in $j$ at fix $k$ comes from the best sequence ending in some $i$ at fix
$k-1$, so
\begin{equation}\label{eq:viterbi}
  \delta_1(j)=\nu_j\,b_j(x_1),\qquad
  \delta_k(j)=\max_i\bigl[\delta_{k-1}(i)\,a_{ij}\bigr]\,b_j(x_k),\qquad
  \psi_k(j)=\arg\max_i\bigl[\delta_{k-1}(i)\,a_{ij}\bigr].
\end{equation}
At the last fix one takes the best phase and follows the pointers $\psi$ backwards~\cite{viterbi1967,forney1973}.
The cost is $K\times S^2$ operations for $S$ phases, against the $S^K$ sequences one would
otherwise compare. The scores shrink at every fix and fall below the smallest number a computer
holds after a few hundred fixes; dividing the column by its sum at each fix, as Vilpellet's code
does, or working with logarithms, keeps them finite and leaves the path unchanged.

\paragraph{The parameters (Baum--Welch).} When $A$, $\nu$ and $b$ are unknown they are fitted by
expectation--maximisation~\cite{baum1970,rabiner1989}. Replacing the maximum of~\eqref{eq:viterbi}
by a sum, a forward and a backward pass give, under the current parameters, the probability
$\gamma_k(j)$ of each phase at each fix and the probability $\xi_k(i,j)$ of each change between
consecutive fixes, given all the symbols. Each round then re-estimates the parameters as expected
counts:
\begin{equation}\label{eq:baumwelch}
  \nu_j=\gamma_1(j),\qquad
  a_{ij}=\frac{\sum_k\xi_k(i,j)}{\sum_k\gamma_k(i)},\qquad
  b_j(x)=\frac{\sum_{k:\,x_k=x}\gamma_k(j)}{\sum_k\gamma_k(j)} .
\end{equation}
No round lowers the probability of the data, and the rounds stop at a local maximum, so the
starting guess matters. The middle formula is the update of $A$ that Vilpellet's training code
leaves out (Section~\ref{sec:conflicts}). The names of the states are not learnt: a state is
``climb'' only because its emissions look like a climb.

\begin{inwords}
Viterbi finds the best phase sequence; Baum--Welch fits the parameters by alternating ``guess the
phases softly'' and ``count''.
\end{inwords}

\subsection{What the model assumes}\label{sec:assumptions}

Three assumptions: phases have no memory, so their durations are geometric; symbols are
independent once the phases are known; the same parameters hold at every fix of every flight.
Vilpellet's symbols are computed over 30-s windows, so two consecutive symbols share 29 of their
30 fixes and are far from independent: the model counts the same evidence about thirty times. The
paper accepts this as approximately true at the scale of 30~s~\cite{vilpellet2026}. In practice
the windows carry the persistence and the HMM settles the changes of phase. None of the
probabilities the model computes is calibrated; only the labels are used.
